
\section{Dämmschichttrockung}
Der Druck in der Prozessluft ist beim Saugen geringer als der Durck im Raum. Der Sättigungsdampfdruck von Wasser beträgt bei einer normalen Raumtemperatur von 20 °C ca. 2.340 Pascal. Damit kann das Wasser verdampfen. 

Der Gesamtdruck in der Dämmschicht sinkt durch den Sauger leicht ab (z. B. von 1013 mbar auf ca. 950 mbar absolut).

\textbf{Frage} Haben wir eine Chance, den Druck am Sensor in der Prozessluft zu messen?

\begin{itemize}
    \item \textbf{Der Effekt auf die Feuchte:} Weil der Druck leicht sinkt, dehnt sich die Luft aus. Ihre Dichte sinkt ein wenig (sie wird geringer als die 1,204 kg/m³).
    \item \textbf{Das Ergebnis in der Praxis:} Die Luft kann durch den veränderten Druck pro Kubikmeter sogar geringfügig mehr Wasserdampf aufnehmen, und die Verdunstung des Wassers in der Dämmschicht wird beschleunigt.
    \item \textbf{Für die Messung:} Die unkorrigierte Magnus-Formel liefert  immer noch einen praxistauglichen Näherungswert. 
\end{itemize}

\textbf{ Diese Aussage sollte hinterfragt werden.}
